\chapter{エントロピー正則化}
\label{ch:seminar-03-entropic}

エントロピー正則化は，
最適輸送問題にエントロピー項を加えることで問題を狭義凸化する手法である．

以下，
\[
  \mathbf{a} \in \R_{++}^{n},\qquad
  \mathbf{b} \in \R_{++}^{m},\qquad
  \sum_{i=1}^{n} a_i = \sum_{j=1}^{m} b_j = 1,
\]
および $\mathbf{C} \in \R_+^{n \times m}$ を固定する．
離散カップリング集合 $\CouplingsD(\mathbf{a}, \mathbf{b})$ は
命題~\ref{prop:sem-discrete-kantorovich} の通り
\[
  \CouplingsD(\mathbf{a}, \mathbf{b})
  =
  \left\{
    \mathbf{P} \in \R_+^{n \times m}
    \;\middle|\;
    \mathbf{P}\ones_m = \mathbf{a},\;
    \mathbf{P}^\top \ones_n = \mathbf{b}
  \right\}
\]
である．

%% ============================================================
%% 4.1 エントロピー正則化
%% ============================================================
\section{エントロピー正則化}
\label{sec:sem-entropic-regularization}

\begin{definition}{離散エントロピー}{sem-discrete-entropy}
  非負行列 $\mathbf{P} \in \R_+^{n \times m}$ に対して，
  \textbf{離散エントロピー}を
  \[
    \Hb(\mathbf{P})
    \defeq
    - \sum_{i=1}^{n} \sum_{j=1}^{m}
      P_{i,j}\bigl(\log P_{i,j} - 1\bigr)
  \]
  で定める．ただし $0\log 0 = 0$ と約束する．
\end{definition}

\begin{remark}{通常の Shannon エントロピーとの違い}{sem-entropy-convention}
  確率行列では $\sum_{i,j}P_{i,j}=1$ なので，
  \[
    \Hb(\mathbf{P})
    =
    -\sum_{i,j} P_{i,j}\log P_{i,j} + 1
  \]
  であり，通常の Shannon エントロピーとは定数 $1$ だけ異なる．
  この定数は最適化の解には影響しない．この形を使う理由は，
  $-\Hb$ の微分が簡単になるからである．積の微分則より
  \[
    \frac{\partial(-\Hb)}{\partial P_{i,j}}
    = \frac{\partial}{\partial P_{i,j}}\bigl[P_{i,j}(\log P_{i,j} - 1)\bigr]
    = (\log P_{i,j} - 1) + P_{i,j} \cdot \frac{1}{P_{i,j}}
    = \log P_{i,j},
  \]
  定義中の $-1$ と積の微分で生じる $+1$ が打ち消し合い，$\log P_{i,j}$ のみが残る．
\end{remark}

\begin{definition}{エントロピー正則化された離散最適輸送}{sem-entropic-ot}
  正則化パラメータ $\varepsilon > 0$ に対して，
  \textbf{エントロピー正則化された離散 Kantorovich 問題}を
  \[
    \MKD_{\mathbf{C}}^\varepsilon(\mathbf{a}, \mathbf{b})
    \defeq
    \min_{\mathbf{P} \in \CouplingsD(\mathbf{a}, \mathbf{b})}
    \left\{
      \inner{\mathbf{C}}{\mathbf{P}}
      - \varepsilon \Hb(\mathbf{P})
    \right\}
  \]
  と定める．この問題の最適解を $\mathbf{P}_\varepsilon$ と書く．
\end{definition}

\begin{proposition}{正則化問題の解の存在と一意性}{sem-entropic-existence-unique}
  任意の $\varepsilon > 0$ に対して，
  正則化問題~\ref{def:sem-entropic-ot} は一意な最適解
  $\mathbf{P}_\varepsilon \in \CouplingsD(\mathbf{a}, \mathbf{b})$ を持つ．
\end{proposition}

\begin{proof}
  $\CouplingsD(\mathbf{a}, \mathbf{b})$ は空でないコンパクト集合である
  （主張~\ref{clm:sem-discrete-kantorovich-existence} の証明）．

  \textbf{目的関数の連続性．}
  $g(x) \defeq x\log x - x$（$x > 0$），$g(0) \defeq 0$ とおく．
  $x > 0$ では $g$ は連続．$x = 0$ では $x = e^{-t}$（$t \to +\infty$）と置換すると
  \[
    x\log x = -\frac{t}{e^t} \to 0,
  \]
  よって $\lim_{x\to 0+} g(x) = 0 = g(0)$ であり $g$ は $[0,\infty)$ 上で連続．
  目的関数 $\mathbf{P} \mapsto \inner{\mathbf{C}}{\mathbf{P}} + \varepsilon\sum_{i,j}g(P_{i,j})$
  は各成分の連続関数の和だから連続であり，コンパクト集合上の連続関数として最小値を達成する．

  \textbf{目的関数の狭義凸性．}
  任意の $x_1, x_2 > 0$，$x_1 \neq x_2$，$\lambda \in (0,1)$ をとり
  $x^* := \lambda x_1 + (1-\lambda)x_2$ とおく．
  平均値定理（定理~\ref{thm:sem-mvt}）を $[x_1, x^*]$ と $[x^*, x_2]$ にそれぞれ適用すると，
  $c_1 \in (x_1, x^*)$，$c_2 \in (x^*, x_2)$ が存在して
  \begin{align*}
    g(x_1) - g(x^*) &= g'(c_1)(x_1 - x^*)
                     = -g'(c_1)(1-\lambda)(x_2 - x_1), \\
    g(x_2) - g(x^*) &= g'(c_2)(x_2 - x^*)
                     = \phantom{-}g'(c_2)\,\lambda(x_2 - x_1).
  \end{align*}
  それぞれ $\lambda$，$(1-\lambda)$ 倍して加えると
  \[
    \lambda g(x_1) + (1-\lambda)g(x_2) - g(x^*)
    = \lambda(1-\lambda)(x_2 - x_1)\bigl[g'(c_2) - g'(c_1)\bigr].
  \]
  $c_1 < c_2$ より $g'(c_1) = \log c_1 < \log c_2 = g'(c_2)$ だから $g'(c_2) - g'(c_1) > 0$，
  また $\lambda(1-\lambda)(x_2-x_1) > 0$ だから右辺 $> 0$．
  よって $g(x^*) < \lambda g(x_1) + (1-\lambda)g(x_2)$ となり，
  $g$ は $(0,\infty)$ 上で狭義凸（定義~\ref{def:sem-convex-function}）．
  $x_1 = 0$，$x_2 > 0$ の場合は $g(0) = 0$ を用いると
  \begin{align*}
    \lambda g(0) + (1-\lambda)g(x_2) - g((1-\lambda)x_2)
    &= (1-\lambda)(x_2\log x_2 - x_2) - \bigl[(1-\lambda)x_2\log(1-\lambda)x_2 - (1-\lambda)x_2\bigr] \\
    &= (1-\lambda)x_2\bigl[\log x_2 - \log(1-\lambda)x_2\bigr] \\
    &= (1-\lambda)x_2\log\frac{1}{1-\lambda} > 0.
  \end{align*}
  よって $g$ は $[0,\infty)$ 上でも狭義凸である．
  各成分の狭義凸関数の和として，目的関数 $\mathbf{P} \mapsto \inner{\mathbf{C}}{\mathbf{P}} + \varepsilon\sum_{i,j}g(P_{i,j})$ も $\mathbf{P}$ について狭義凸である．

  \textbf{最適解の一意性．}
  目的関数が狭義凸であるとき，最適解は高々1つである．
  実際，$\mathbf{P}_1 \neq \mathbf{P}_2$ がともに最適解であるとすると，
  中点 $\mathbf{P}^* = (\mathbf{P}_1 + \mathbf{P}_2)/2$ も $\CouplingsD(\mathbf{a},\mathbf{b})$ に属し（凸集合），
  狭義凸性から
  \[
    f(\mathbf{P}^*) < \tfrac{1}{2}f(\mathbf{P}_1) + \tfrac{1}{2}f(\mathbf{P}_2) = f(\mathbf{P}_1)
  \]
  となり，$\mathbf{P}_1$ が最適解であることに矛盾する．よって最適解は一意である．
\end{proof}

\begin{proposition}{正則化解の正値性}{sem-entropic-positive-solution}
  正則化問題の一意解 $\mathbf{P}_\varepsilon$ は
  \[
    (P_\varepsilon)_{i,j} > 0
    \qquad(\forall\, i \in \range{n},\, j \in \range{m})
  \]
  を満たす．
\end{proposition}

\begin{proof}
  背理法で示す．ある $(i,j)$ で
  $(P_\varepsilon)_{i,j}=0$ とする．
  $a_i>0$ なので同じ行に $(P_\varepsilon)_{i,j_1}>0$ となる
  $j_1$ が存在する．また $b_j>0$ なので同じ列に
  $(P_\varepsilon)_{i_1,j}>0$ となる $i_1$ が存在する．
  小さい $\theta>0$ に対し，4成分だけを
  \[
    P_{i,j} \leftarrow P_{i,j}+\theta,\quad
    P_{i,j_1} \leftarrow P_{i,j_1}-\theta,\quad
    P_{i_1,j} \leftarrow P_{i_1,j}-\theta,\quad
    P_{i_1,j_1} \leftarrow P_{i_1,j_1}+\theta
  \]
  と変化させる．$\theta$ を
  $(P_\varepsilon)_{i,j_1}$ と $(P_\varepsilon)_{i_1,j}$ より小さく取れば，
  非負性と周辺条件は保たれる．

  この変化による線形コストの変化は $\theta$ に比例する．
  一方，エントロピー項のうち $0$ だった成分に対応する負エントロピー
  $x\log x-x$ の変化は
  \[
    \theta\log\theta - \theta
  \]
  であり，これを $\theta$ で割ると $\log\theta-1 \to -\infty$
  となる．したがって十分小さい $\theta>0$ では，
  全体の目的関数
  $\inner{\mathbf{C}}{\mathbf{P}}-\varepsilon\Hb(\mathbf{P})$
  は小さくなる．これは $\mathbf{P}_\varepsilon$ の最適性に矛盾する．
  よって全成分が正である．
\end{proof}

\begin{remark}{$\varepsilon$ の役割}{sem-epsilon-role}
  $\varepsilon$ は，元の線形計画と独立カップリングの間を補間する量である：
  \begin{itemize}
    \item $\varepsilon \to 0$ のとき，$\mathbf{P}_\varepsilon$ は
      非正則化問題 $\MKD_{\mathbf{C}}$ の最適解のうち
      エントロピーが最大のものに収束する
      （定理~\ref{thm:sem-entropic-convergence-eps-zero}）．
    \item $\varepsilon \to +\infty$ のとき，$\mathbf{P}_\varepsilon$ は
      積測度 $\mathbf{a} \mathbf{b}^\top$ に収束する
      （注意~\ref{rem:sem-entropic-convergence-eps-infty}）．
  \end{itemize}
\end{remark}

\subsection{KL ダイバージェンスによる定式化}
\label{subsec:sem-kl-formulation}

\begin{definition}{離散 KL ダイバージェンス}{sem-kl-divergence}
  非負行列 $\mathbf{P}, \mathbf{K} \in \R_+^{n \times m}$ に対して，
  \textbf{KL ダイバージェンス}を
  \[
    \KLD(\mathbf{P} \| \mathbf{K})
    \defeq
    \sum_{i,j}
    \left(
      P_{i,j}\log\frac{P_{i,j}}{K_{i,j}}
      - P_{i,j} + K_{i,j}
    \right)
  \]
  で定める．ただし $0\log 0 = 0$ と約束し，
  $P_{i,j} > 0$ かつ $K_{i,j}=0$ となる成分がある場合は
  $\KLD(\mathbf{P}\|\mathbf{K}) = +\infty$ とする．
\end{definition}

\begin{definition}{Gibbs カーネル}{sem-gibbs-kernel}
  コスト行列 $\mathbf{C}$ と $\varepsilon > 0$ に対して，
  \textbf{Gibbs カーネル} $\mathbf{K} \in \R_{++}^{n \times m}$ を
  \[
    K_{i,j} \defeq \exp\!\left(-\frac{C_{i,j}}{\varepsilon}\right)
  \]
  で定める．
\end{definition}

\begin{proposition}{正則化 OT は KL 射影である}{sem-entropic-kl-projection}
  $\mathbf{K} = \exp(-\mathbf{C}/\varepsilon)$ とすると，
  正則化問題~\ref{def:sem-entropic-ot} の最適解は
  \[
    \mathbf{P}_\varepsilon
    =
    \argmin_{\mathbf{P} \in \CouplingsD(\mathbf{a}, \mathbf{b})}
    \KLD(\mathbf{P}\|\mathbf{K})
  \]
  と書ける．
\end{proposition}

\begin{proof}
  $\log K_{i,j} = -C_{i,j}/\varepsilon$ より
  \begin{align*}
    \varepsilon \KLD(\mathbf{P}\|\mathbf{K})
    &=
    \varepsilon\sum_{i,j}
    \left(
      P_{i,j}\log P_{i,j}
      - P_{i,j}\log K_{i,j}
      - P_{i,j}
      + K_{i,j}
    \right) \\
    &=
    \sum_{i,j} C_{i,j}P_{i,j}
    + \varepsilon\sum_{i,j} P_{i,j}(\log P_{i,j}-1)
    + \varepsilon\sum_{i,j} K_{i,j} \\
    &=
    \inner{\mathbf{C}}{\mathbf{P}} - \varepsilon \Hb(\mathbf{P})
    + \varepsilon\sum_{i,j} K_{i,j}.
  \end{align*}
  最後の項は $\mathbf{P}$ に依存しない定数なので，
  両者の最小化問題は同じ最適解を持つ．
\end{proof}

\subsection{正則化パラメータの極限}
\label{subsec:sem-epsilon-limits}

\begin{theorem}{$\varepsilon \to 0$ による非正則化 OT への収束}{sem-entropic-convergence-eps-zero}
  $\varepsilon_k > 0$，$\varepsilon_k \to 0$ とし，
  $\mathbf{P}_{\varepsilon_k}$ を正則化問題の一意解とする．
  もし部分列 $\mathbf{P}_{\varepsilon_{k_\ell}}$ が
  $\mathbf{P}^0$ に収束するならば，
  $\mathbf{P}^0$ は非正則化問題
  \[
    \min_{\mathbf{P} \in \CouplingsD(\mathbf{a}, \mathbf{b})}
    \inner{\mathbf{C}}{\mathbf{P}}
  \]
  の最適解である．また最適値は
  \[
    \lim_{\varepsilon \to 0}
    \MKD_{\mathbf{C}}^\varepsilon(\mathbf{a},\mathbf{b})
    =
    \MKD_{\mathbf{C}}(\mathbf{a},\mathbf{b})
  \]
  を満たす．さらに，非正則化問題の最適解が複数ある場合，
  $\mathbf{P}_\varepsilon$ はその中でエントロピー $\Hb$ が最大のものに収束する．
\end{theorem}

\begin{proof}
  非正則化問題の任意の最適解を $\mathbf{P}^*$ とする．
  $\mathbf{P}_\varepsilon$ の最適性より
  \[
    \inner{\mathbf{C}}{\mathbf{P}_\varepsilon}
    - \varepsilon\Hb(\mathbf{P}_\varepsilon)
    \leq
    \inner{\mathbf{C}}{\mathbf{P}^*}
    - \varepsilon\Hb(\mathbf{P}^*).
  \]
  一方，$\mathbf{P}^*$ は非正則化問題の最適解なので
  \[
    0
    \leq
    \inner{\mathbf{C}}{\mathbf{P}_\varepsilon}
    - \inner{\mathbf{C}}{\mathbf{P}^*}
    \leq
    \varepsilon\bigl(\Hb(\mathbf{P}_\varepsilon)-\Hb(\mathbf{P}^*)\bigr).
  \]
  $\CouplingsD(\mathbf{a}, \mathbf{b})$ はコンパクトであり，$\Hb$ はその上で有界なので，
  右辺は $\varepsilon \to 0$ で $0$ に収束する．
  したがって
  $\inner{\mathbf{C}}{\mathbf{P}_\varepsilon}
  \to \inner{\mathbf{C}}{\mathbf{P}^*}$ である．また
  $\varepsilon\Hb(\mathbf{P}_\varepsilon) \to 0$ なので，
  $\MKD_{\mathbf{C}}^\varepsilon(\mathbf{a},\mathbf{b})
  \to \MKD_{\mathbf{C}}(\mathbf{a},\mathbf{b})$ である．
  任意の極限点 $\mathbf{P}^0$ は非正則化問題の最適値を達成する．
  さらに上の不等式から
  \[
    \Hb(\mathbf{P}_\varepsilon) \geq \Hb(\mathbf{P}^*)
  \]
  が任意の非正則化最適解 $\mathbf{P}^*$ について成り立つ．
  $\Hb$ は連続なので，任意の極限点 $\mathbf{P}^0$ は
  非正則化最適解の中で $\Hb$ を最大化する．
  非正則化最適解集合は凸であり，$-\Hb$ は狭義凸であるから，
  $\Hb$ を最大化する非正則化最適解は一意である．したがって
  $\mathbf{P}_\varepsilon$ 全体がその最大エントロピー最適解に収束する．
\end{proof}

\begin{remark}{$\varepsilon \to +\infty$ による独立カップリングへの収束}{sem-entropic-convergence-eps-infty}
  $\varepsilon \to +\infty$ のとき
  $\mathbf{P}_\varepsilon \to \mathbf{a}\mathbf{b}^\top$ が成り立つ．
  $\varepsilon$ で目的関数を割ると
  $\varepsilon^{-1}\inner{\mathbf{C}}{\mathbf{P}} - \Hb(\mathbf{P})$
  の最小化であり，$\varepsilon \to +\infty$ では
  エントロピー最大化
  $\max_{\mathbf{P} \in \CouplingsD} \Hb(\mathbf{P})$
  に帰着する．周辺分布を固定したエントロピー最大化の解は
  独立カップリング $\mathbf{a}\mathbf{b}^\top$ である．
\end{remark}

